\begin{abstract}
Performance in modern GPU-centric systems is increasingly defined by resource management policies, such as memory placement, scheduling, and observability, rather than just raw compute power. Currently, these policies are implemented at suboptimal abstraction layers: either hard-coded into closed, static drivers and firmware, or siloed within user-space runtimes that lack transparency, cross-tenant visibility, and fine-grained hardware control. We argue that the OS GPU driver and device layer must serve as the system's programmable narrow waist. However, extending standard host-side eBPF is insufficient: it remains blind to critical device-side events, and directly injecting policy code into GPU kernels risks compromising safety and efficiency.

We present \sys, a unified policy plane that treats the GPU driver and device as a single programmable OS subsystem. \sys extends GPU drivers to expose safe hooks and introduces a device-side eBPF runtime that executes verified policy logic efficiently within GPU kernels, enabling coherent, application-transparent policies that coordinate across the host-device boundary. Evaluation on realistic inference and training workloads shows that \sys improves throughput by up to  5× and reduces tail latency by 2× with negligible overhead, all without modifying applications or restarting drivers.
% \xiangyu{I think it is better to reduce the abstract content by 50-60\% by removing some example details.}
% \yusheng{should we use general "policy" or make it resources management policy? we are also doing  optimization for a single application and doing observability/tracing}
% \JC{
% Resource management policies have become crucial in modern GPU systems, as diverse workloads contend for limited resources in multi-tenant environments. However, current approaches to policy specification and enforcement are inadequate. Policies are either hard-coded into opaque drivers or scattered across user-space runtimes, lacking flexibility, coordination, and strong safety guarantees.

% We propose exposing the GPU driver and device layer as a programmable resource management substrate using eBPF, enabling application-transparent policies with stack-wide coordination, full hardware control, and strong safety guarantees. \sys introduces a verifiable runtime that executes management policies within the GPU stack, ensuring policy code is checked for bounded execution and memory access. This allows policies to dynamically adapt to workload requirements while guaranteeing system stability. Our evaluation on real-world workloads shows that \sys improves


% }
\end{abstract}

\section{Introduction}

% 1. (General problem) GPU is important (1). Motivate the resources management policy is important (resource mpolicy has emerged as ... 
% 2. (The problem with current works) Current design rigiet for us to create new resource management policy, some one has try to modify drivers; others are trying to do it in usersapce framwork
% 3. Our key insight is that the GPU driver is uniquely positiooned in the GPU stack to have transparency into all applications, etc. 
% 4. (Our claim) In this paper, we present \sys{}, a new system that enables flexible resource management policies for GPU deployments.  Building on our key insight, \sys{} implements a stable set of control plane hooks that allow for flexible resource managment policies...
% 5. To support flexibile policies, applications must be able to observe an application's behavior on the GPU device. Safety is required... Thus, we implement our SIMT verifier. 
%6. Finally, to allow the driver and device components to communicate, we need a communication substrate.
%7. We implemented \sys{} on.... We evaluated on xyz workloads and our results are awesome. 
%8. Our contributions are as follows: 
%   a. 


% 1. (General problem) GPU is important (1). Motivate the resources management policy is important (resource mpolicy has emerged as ...
GPUs account for an increasingly large share of compute resources and capital expenditures in modern systems~\cite{jeon2019analysis}, supporting diverse workloads such as latency-sensitive inference~\cite{kwon2023efficient,agrawal2024taming,yu2024orca,}, compute intensive training~\cite{shoeybi2019megatron,rasley2020deepspeed},  graph analytics~\cite{lin2024towards,chen2020pangolin,wang2016gunrock}, vector search~\cite{pan2024survey,cao2023gpu}, and data visualization~\cite{moreland2016vtk}. To support such workloads, modern GPU software stacks typically consist of user-space frameworks and runtimes, OS kernel drivers, vendor firmware, and device-side GPU kernel code executing directly on GPU hardware. Given this multi-layered, heterogeneous architecture spanning compute, memory, and software boundaries, effective performance and efficiency depend not only on raw compute power, but critically on adaptive resource management \emph{policies}. Such policies orchestrate data placement and migration between host and device memory~\cite{sun2024llumnix,zhong2024distserve,patel2024splitwise,strati2024orion,sheng2023flexgen}, manage kernel admission, scheduling, and preemption in multi-tenant deployments~\cite{fu2024serverlessllm,shubha2024usher,hu2023lucid,li2022miso,jayaram2023sia,xiao2018gandiva,narayanan2020heterogeneity}, and enable dynamic observability for online adaptation and optimization~\cite{agrawal2024vidur}. Crucially, these policies must continuously evolve alongside changing hardware architectures, emerging workload characteristics, and varying service-level agreements (SLAs)~\cite{kwon2023efficient,yu2020fine,delimitrou2013paragon,delimitrou2014quasar,lo2015heracles,qiao2021pollux}.\arq{The way these are cited, it appears as though all of the papers apply to the varying SLAs part of the sentence. Do they?} \yusheng{No, we should fix that}
\arq{This paragraph is trying to say too much, and all the details are making it difficult for the reader to glean the main idea. Every paragraph should boil down to one idea, which should be conveyed in a single topic sentence.  The easiest way to write is to have the topic sentence be the first sentence in the paragraph; the thing that comes first.}
\arq{I think that the main point of this paragraph, and thus my first sentence, would be: ``The performance and efficiency of GPU-based systems increasingly depends not only on raw compute and memory performance, but also on resource management policies.''}
\yusheng{One thing I'm wondering is that whether we should use one sentence to introduce the background of GPU stacks.}\arq{We shouldn't be looking at "fixes" that involve making this paragraph larger---there's already too much going on in this one place.  I think there are two options: First, try to write the paragraph by starting from the suggestion that I had of "policies matter". Then, we would explain why policies (i.e., because of current GPU stacks) afterwards. Second, we could create two paragraphs, one that says that GPU stacks are complex but important and the second that says that policies are super important. The slightly *easier writing approach is the second approach that I advocated for, but I think the first one would probably actually be a better paper}
% enhance the multi-tenant


% 2. (The problem with current works) Current design rigiet for us to create new resource management policy, some one has try to modify drivers; others are trying to do it in usersapce framwork
Despite their importance, current GPU resource management policies remain constrained by rigid system abstractions. Vendor GPU drivers and firmware enforce static, opaque mechanisms lacking programmability and adaptability~\cite{nvidia-open-gpu}, limiting performance optimization. Academic works like TimeGraph and Gdev shift GPU scheduling and memory control into Linux device drivers~\cite{kato2011timegraph,kato2012gdev}, and recent systems refactor driver-level control paths~\cite{kernelintercept,fan2025gpreempt,wang2024gcaps}, but implement policies in static kernel modules. Lacking safe programmable interfaces, adapting policies requires driver changes, causing operational and safety issues. To achieve flexibility without kernel modifications, recent user-space serving systems retrofit programmability above drivers: Paella~\cite{ng2023paella} and XSched~\cite{shen2025xsched} expose software-defined scheduling, while Pie~\cite{gim2025pie} and KTransformers~\cite{chen2025ktransformers} provide programmable runtimes for LLM inference. 
However, user-space approaches face three fundamental limitations. First, programmability is application-bound~\cite{ng2023paella,wang2024gcaps}: policies couple to specific runtimes and cannot govern arbitrary, unmodified applications. Deploying new policies requires invasive code changes or porting applications to specific stacks, limiting reuse. Second, user-space runtimes lack global coordination. Isolated in process silos, they cannot safely share memory or bandwidth across tenants, forcing conservative resource pre-allocation, preventing dynamic sharing, and lowering utilization~\cite{wu2024streambox}. Third, user-space systems lack fine-grained execution control. Operating above drivers, they cannot act at hardware-interrupt or device-context granularity, limiting precise preemption and compute–transfer overlap~\cite{fan2025gpreempt}. Profiling-oriented instrumentation frameworks~\cite{cupti,villa2019nvbit,neutrino,egpu} profile GPU kernels but remain isolated from host-side drivers, lacking programmability and cross-layer coordination.\arq{There are too many ideas in this paragraph without an overarching theme. You want to say at the top of the paragraph what the main point of all of the things will be.  I don't think the answer is "current GPU resource management policies remain constrained by rigid abstraction"---what is a rigid abstraction? Are the userspace approaches or the kernel module approaches rigid? I don't think that's an accurate and all-encompassing complaint about the whole body of work.  We really want to summarize the issues with the prior work to get to the crux of why they fail.} \yusheng{I don't know what is igid abstraction...I heard it from you yesterday, we can come up with a better word. I think so, we should have a more clear expression.}\arq{I think that it is a tradeoff: ``current state of the art approaches force developers to express GPU resource management policies that are either ineffective or unsafe.  We argue that the userspace programs are ineffective, since they cannot get visibility and management of low-level GPU stack behaviors. Then, we argue that kernel approachs that involve custom drivers and kernel modules are inherently unsafe.}

\yusheng{Maybe the key idea we want to express si that there is no policy inside the OS layer?}

% 3. Our key insight is that the GPU driver is uniquely positiooned in the GPU stack to have transparency into all applications, etc.
    To remedy the above limitations, we argue that effective GPU resource management requires an \textbf{OS-level narrow waist for policy}. \arq{the english in this phrase does not make sense. It also is not obvious why a narrow waist for supporting resource management policies would be helpful for the issues at stake, although that's in part because the previous paragraph is not very clear.} Specifically, the GPU driver layer is uniquely positioned: it offers application transparency (policies can be adapted without modifying applications), global cross-tenant resource coordination, and direct control over privileged mechanisms such as page migration, kernel scheduling, preemption, and memory prefetching. Recent CPU-side frameworks, such as eBPF, have demonstrated the feasibility of dynamic, safe, kernel-level scheduling and resource management policies~\cite{bachl2021flow,Zhong22,jia2023programmable,schedext-docs,zussman2025cache_ext}. Inspired by this success, we propose treating GPU resource management as a first-class, programmable kernel subsystem. However, host-only eBPF is insufficient because critical scheduling decisions, such as thread-block dispatch, warp execution, and device-side memory accesses, occur directly on GPU hardware and remain invisible and uncontrollable to the host.
A practical solution requires exposing programmable policy hooks within GPU drivers and executing policy logic inside GPU kernels safely and efficiently.\arq{Again, this paragraph is too many ideas packed into one. What is the main topic sentence? I'm not sure what the narrow waist thing is even trying to convey if I'm honest.  My intuition is that it probably isn't actually what we think is novel, since prior work has suggested the idea of implementing policies in the GPU driver before.} \yusheng{Prev work does not talk about programmable policies in GPU driver. alrought there are some mechaism and policy changes to improve it.}\arq{The kernel drivers that we cited were policies that were implemented by modifying/creating a new GPU driver. So, that's what I mean: people know that you can/should/might want to do policies at the driver level because the driver is "the OS" essentially.}\arq{I think the core novelty is more about the simple resource management policy interface that is at once low enough level to have visibility and manage low-level details but high-level enough to simplify development and enable safe programming approaches}

% 4. (Our claim) In this paper, we present \sys{}, a new system that enables flexible resource management policies for GPU deployments.  Building on our key insight, \sys{} implements a stable set of control plane hooks that allow for flexible resource managment policies...
% 5. To support flexibile policies, applications must be able to observe an application's behavior on the GPU device. Safety is required... Thus, we implement our SIMT verifier. 
%6. Finally, to allow the driver and device components to communicate, we need a communication substrate.
We present \sys{}, a cross-layer policy runtime that transforms the GPU driver and device into a programmable OS subsystem. \sys{} exposes stable, verified control-plane hooks within the GPU driver, enabling safe and dynamically programmable resource management policies as eBPF programs. \sys{} further supports injecting verified eBPF policy logic directly into GPU kernels, providing consistent and transparent policy enforcement across the host-device boundary.
However, designing \sys{} posed three primary technical challenges, which we address as follows: \arq{Stylistically, I don't like how we label these with C1, C2, and C3 and then go into the explanation. I just don't care for the style.}

\textbf{C1. Safe and Expressive driver extensibility.} Current GPU drivers were never designed to expose a programmable policy interface: directly surfacing low-level mechanisms (e.g., page tables, interrupts, runlists) risks kernel crashes and driver instability, while overly abstract interfaces sacrifice the expressiveness needed for sophisticated resource management. To address this tension, we introduce a narrow, stable GPU policy interface built around safe, hardware-aligned abstractions (regions, queues), dynamically separating policy decisions from low-level driver mechanisms to provide rich programmability without compromising system safety.

%\xiangyu{The last sentence is a little bit abstract for me to understand. Is it possible to concretize it a little bit?} \arq{It isn't a challenge if we solved it by just doing it.  Technical challenges are things that require novelty to resole. Was there some novelty to resolve this challenge? If not, then it isn't a technical challenge and we need to remove it.  We had this conversation yesterday, no?} \yusheng{yes, we chat yesterday, I think there are some here and do not want to remove this point, just I want to make sure we can express it clearly. I think the novelty is that, if everybody knows how to do, and you just do it, there is no novelty; but seems people don't have a clear idea how to do it. I think we can just say with the similar approach of how cacheext say their idea.}\arq{Yusheng, this is not a challenge as presented. } \yusheng{I think the key things is that we want to presented it clearly? Any suggestions? like cacheext?} \arq{My suggestion is to delete the challenge because it isn't a challenge.} \yusheng{Do you think there is no novelty in cachext and no challenge?}\arq{that is a false analogy. This challenge says "it is difficult to identify where to build our interface because the driver is tightly coupled". And, we solved that challenge by.... identifying where to build our interface.  Challenges cannot be solved by "we worked hard". That isn't a viable technical challenge for a research paper.}
% \yusheng{I agree, maybe the coupled is not a real challenge, but i think we do have some design challnge in driver part, like cacheext? How do they state?}\arq{As we discussed yesteray, Cacheext challenges were safety, scalability, etc. They were adjectives that they argued were important for the system rather than specific tasks that needed to be done. }
% \yusheng{I agree, we should rewrite this into something like how we make the driver interface safe and scalable, and more like how they did it?}
% \arq{We should write it in the way that we feel best expresses the work that we did, just copying cacheext is not a good approach.  One way to do it might be to discuss the high level adjectives that we want (scalability, safety, efficiency, etc.).  But, in my mind, I think it would be reasonable for this particular challenge to disappear and discuss our approach to this problem earlier in the text (like, when we discuss our key insight, since we are already presenting it there). Then, the challenges would be Safe+efficient driver-side execution and unified state management.  It is OK to just have those two.  My gut is that the best version of our paper would go with the high-level adjectives rather than the low level challenges that I've described. So, that is probably *the right thing*. But, I am worried about whether we'll actually be able to deliver on those writing changes, esp. since I cannot spend much of tomorrow on this. }
% \yusheng{One thing I want to mention is that if we remove it or put it as high level, it seems all we claim good/novel is the device part. However, most of the evaluation improvement and our high level statement comes from the driver part...}
% \arq{I think that's reductive to what I'm suggesting. We can claim that the interface is the key novelty, but not that it is a technical innovation required to solve a technical challenge.  In other words, we can say that our key novelty is in designing a programmable interface for resource management policies, but we don't have to say "it was so hard" like how this is written currently.  I think the text about the interface should come earlier in the introduction, up where we describe our key insight.}
% \yusheng{I think we might want some like this, but 1-2 points and must high level?}
% \yusheng{cachebpf’s design is motivated by four main insights. First, modern storage devices are very fast and support millions of IOPS, so custom page cache policies must run with low overhead. Therefore, we design cachebpf so that its eBPF-based policies run in the kernel, avoiding expensive and frequent synchronization between the kernel and userspace. Second, caching algorithms are very diverse and may use complex data structures. To address this challenge, cachebpf exposes a simple yet flexible interface that allows applications to define one or more variable-sized lists of pages, and a set of policy functions (e.g., admission, eviction) that operate on these lists, which can be used to express a wide range of eviction policies. Third, in order for cachebpf to be useful in multi-tenant scenarios, it should allow each application to use its own policy without interfering with others. We identify cgroups as a natural isolation boundary. Thus, cachebpf allows each cgroup to implement its own eviction policy without interfering with other cgroups. Finally, custom policies determine which pages to evict and return page references to the kernel. However, these references may be invalid, which could lead to kernel crashes or security breaches. To solve this, cachebpf maintains a registry of valid page references, which is used to validate the page references returned by the user-defined policies.}

\textbf{C2. Safe and efficient device-side execution.} Executing policy logic efficiently and safely within GPU kernels poses challenges due to the GPU's massively parallel SIMT execution environment. Device-side policy programs must bound resource usage and guarantee safety without impacting critical-path performance. \sys{} resolves this with a specialized SIMT-aware eBPF verifier extension and optimized GPU-side runtime that transparently inject verified policy logic directly into GPU kernels.

\textbf{C3. efficient and consistent state management.} Policies must correlate state across host and device boundaries, a task complicated by their heterogeneous memory hierarchies. \sys{} addresses this challenge by employing cross-device eBPF maps, enabling unified state management within a single control plane.


We implement \sys{} on Linux by extending the NVIDIA open GPU kernel modules~\cite{nvidia-open-gpu} and providing an eBPF runtime for GPU devices. We evaluate \sys across realistic GPU workloads including LLM inference pipelines, GNN training, vector search, and mixed-priority multi-tenant scenarios. Using \sys, we implement adaptive memory prefetching and eviction policies, fine-grained kernel preemption strategies, dynamic work-stealing schedulers, and programmable GPU kernel observability tools. \sys-based policies improve throughput by up to $5\times$ and reduce tail latency by up to $2\times$ compared to static heuristics, while instrumentation-only deployments incur only $2$--$3\%$ overhead. These results demonstrate that unified eBPF-based programmability spanning host and device is a practical abstraction for GPU-centric systems.

Our contributions can be summarized as follows:

\begin{itemize}[leftmargin=*,itemsep=0pt]
  \item We identify key programmability gaps in GPU stacks for memory management, scheduling, and observability, showing how the lack of a programmable, kernel-anchored policy waist at the GPU driver and device layer impedes adaptation to customized requirements. % evolving workloads and SLAs.
  \arq{Where in the paper do we present this novel contribution?} \yusheng{Fig1 and Fig2?}
  \item We design and implement \sys, an eBPF-based policy runtime treating the GPU driver and device execution layer as a OS subsystem. \sys exposes verified driver hooks for memory and scheduling, executes the same eBPF IR inside GPU kernels via a device-side JIT, and shares state through unified-memory-backed BPF maps under a single verifier and control plane.
  \item We demonstrate that \sys enables adaptive, workload-specific policies that improve performance, resource utilization, and tail latency on realistic LLM inference, training, and vector-search workloads.
\end{itemize}
